\section*{Capítulo \ref{ch:b2:alkene-redox} --- Oxidación y reducción de los alquenos}

\begin{solution}{exo:b2:alkene-redox:1}
Metilciclohexano; 1-metilciclohexan-1-ol; trans-2-metilciclohexan-1-ol (adición sin de
H y B, de modo que el \ce{OH} y el nuevo H quedan del mismo lado, y el metilo en trans respecto al \ce{OH});
óxido de 1-metilciclohexeno; cis-1-metilciclohexano-1,2-diol; 6-oxoheptanal
\ce{CH3CO(CH2)4CHO}.
\end{solution}

\begin{solution}{exo:b2:alkene-redox:2}
Hidroboración--oxidación: 2-metilpropan-1-ol. Hidratación catalizada por ácido:
2-metilpropan-2-ol.
\end{solution}

\begin{solution}{exo:b2:alkene-redox:3}
trans-2,3-Dimetiloxirano, con los dos metilos en lados opuestos del anillo. Es quiral (sin
plano de simetría) y se forma como mezcla racémica.
\end{solution}

\begin{solution}{exo:b2:alkene-redox:4}
Tratamiento reductor: dos moléculas de propanal. Tratamiento oxidante: dos moléculas de ácido
propanoico.
\end{solution}

\begin{solution}{exo:b2:alkene-redox:5}
(a) meso-butano-2,3-diol, inactivo por ser aquiral. (b) los dioles
(2\textit{R},3\textit{R}) y (2\textit{S},3\textit{S}) en cantidades iguales, inactivos
por ser racémicos.
\end{solution}

\begin{solution}{exo:b2:alkene-redox:6}
En medio ácido el agua ataca el carbono terciario; en medio básico el hidróxido ataca el \ce{CH2}; en ambos
casos acaba un \ce{OH} en cada carbono: el mismo 2-metilpropano-1,2-diol. Con metanol las dos
vías dan éteres distintos (como en el capítulo).
\end{solution}

\begin{solution}{exo:b2:alkene-redox:7}
Hidrogenación con un equivalente de \ce{H2} sobre un catalizador de paladio envenenado
(Lindlar): la adición sin al triple enlace se detiene en el (\textit{Z})-alqueno.
\end{solution}

\begin{solution}{exo:b2:alkene-redox:8}
La propanona (3 C) y el butanal (4 C), unidos por sus carbonos carbonílicos, dan el 2-metilhex-2-eno,
\ce{(CH3)2C=CH-CH2CH2CH3}.
\end{solution}

\begin{solution}{exo:b2:alkene-redox:9}
\ce{BH3} (o un borano voluminoso), y luego \ce{H2O2} y \ce{NaOH}: hexan-1-ol. La hidratación catalizada por ácido
pasa por el carbocatión secundario y da hexan-2-ol (Markovnikov), con
subproductos de transposición.
\end{solution}

\begin{solution}{exo:b2:alkene-redox:10}
El \omterm{def:b2:alkene-redox:epoxide}{peroxiácido}, un electrófilo, reacciona más deprisa con el doble enlace del anillo, más sustituido y más rico
en electrones. El borano voluminoso reacciona con el \ce{C=CH2} terminal, menos impedido.
\end{solution}

\begin{solution}{exo:b2:alkene-redox:11}
trans: \omterm{def:b2:alkene-redox:epoxide}{peroxiácido}, y luego agua en medio ácido (apertura anti). cis: \ce{OsO4} catalítico con un
reoxidante, o permanganato diluido en frío (sin).
\end{solution}

\begin{solution}{exo:b2:alkene-redox:12}
\ce{C6H10} tiene dos grados de insaturación; un equivalente de \ce{H2} significa un \ce{C=C}, luego
un anillo. Un único dialdehído de seis carbonos procede de un anillo abierto por su doble enlace:
el ciclohexeno.
\end{solution}

\begin{solution}{pb:b2:alkene-redox:1}
\textbf{1.} $10 \times 12.011 + 16 \times 1.008 = \qty{136.24}{g/mol}$.
\textbf{2.} $(2 \times 10 + 2 - 16)/2 = 3$.
\textbf{3.} Un anillo y dos dobles enlaces.
\textbf{4.} \ce{C10H16 + 2H2 -> C10H20}, 1-isopropil-4-metilciclohexano.
\textbf{5.} $1.00/136.24 = \qty{7.34}{mmol}$ de limoneno, \qty{14.68}{mmol} de \ce{H2}.
\textbf{6.} $V = nRT/p = 0.01468 \times 8.314 \times 298.15/10^5 = \qty{3.64e-4}{m^3}$,
\qty{364}{mL}.
\textbf{7.} No: los carbonos 1 y 4 del anillo llevan cada uno dos caminos idénticos por el anillo; los isómeros cis y
trans son aquirales.
\textbf{8.} Metanal \ce{HCHO} (1 C), del extremo \ce{=CH2}, y un único compuesto \ce{C9},
el 3-acetil-6-oxoheptanal \ce{CH3CO-CH2CH2-CH(COCH3)-CH2-CHO}.
\textbf{9.} Sus dos carbonos pertenecen al mismo anillo: romper el doble enlace abre el anillo
sin separar nada.
\textbf{10.} El metanal (un gas).
\textbf{11.} Los grupos aldehído pasan a ácidos carboxílicos (el metanal llega hasta \ce{CO2}); las
cetonas no cambian.
\textbf{12.} La cetona y el aldehído de la cadena \ce{C9} estaban unidos en el anillo; la segunda
cetona (\ce{COCH3}) y el metanal eran los dos extremos del doble enlace de la cadena lateral.
\textbf{13.} Uno (el C4 del anillo); la \omterm{def:b2:alkene-redox:cleavage}{ozonólisis} no lo toca y sobrevive como
estereocentro del producto \ce{C9}.
\textbf{14.} El \ce{C=CH2} exocíclico, menos impedido que el doble enlace trisustituido del anillo.
\textbf{15.} 2-(4-Metilciclohex-3-en-1-il)propan-1-ol: \ce{CH2OH} en el antiguo carbono
terminal.
\textbf{16.} Primario.
\textbf{17.} El \ce{OH} va al carbono menos sustituido, al contrario de la orientación
predicha por la regla de Markovnikov para la hidratación.
\textbf{18.} Dos configuraciones en el nuevo centro, que dan dos diastereoisómeros (con el centro
ya existente).
\textbf{19.} La hidratación catalizada por ácido (u otra hidratación de Markovnikov): el alcohol terciario,
el $\alpha$-terpineol.
\textbf{20.} El doble enlace del anillo: trisustituido, más rico en electrones.
\textbf{21.} 1,2-Epoxi-1-metil-4-(prop-1-en-2-il)ciclohexano, con el oxígeno como puente entre C1 y C2.
\textbf{22.} En C1, el carbono más sustituido (el que lleva el metilo).
\textbf{23.} 1-Metil-4-(prop-1-en-2-il)ciclohexano-1,2-diol, con los dos \ce{OH} en trans (apertura
anti).
\textbf{24.} Una \omterm{def:b2:alkene-redox:dihydroxylation}{dihidroxilación} sin del doble enlace del anillo (\ce{OsO4} catalítico), que da el diol
cis, si se consigue que el reactivo ataque selectivamente ese doble enlace.
\textbf{25.} \qty{1.00}{g} de limoneno absorbe $\boldsymbol{\approx \qty{364}{mL}}$ de
dihidrógeno.
\end{solution}
